\section{A grazing-uniform resolvent for the physical reflection equations}
\label{sec:gr-resolvent}
We retain the actual circular-scatterer trajectories, radius interval,
unit-speed clock, and all exact selectors of \eqref{eq:v76-event}.
The issue here is the derivative of an endpoint boundary-value problem,
not a derivative bound for the expanding forward collision map.
All derivatives are taken at regular physical points of a fixed actual
itinerary. No unvisited disk is removed from its admissible domain.

Put $a_*=R_+/\tau_-=47/6$, $\lambda=50/47$, and $I=\{1,\ldots,m-1\}$.
A sum over $I$ is empty when $m=1$. Write
$X_i=x_{i+1}-x_i$, $\vartheta_i=\arg X_i$, using local real lifts of
these angles. The specular law is equivalent, on the physical branch,
to constancy of
\begin{equation}\label{eq:gr-reflection}
 E_i(\alpha)=2\alpha_i-\vartheta_{i-1}(\alpha)-\vartheta_i(\alpha),
 \qquad i\in I.
\end{equation}
Its constant is an odd multiple of $\pi$. The value of that multiple
does not enter a derivative; it is fixed on each local lift.
The reflection equation is not obtained by assuming that division by
a small incidence is harmless.

\begin{lemma}[The full physical Jacobian]\label{lem:gr-jacobian}
At an actual regular trajectory, set $w_i=\tau_i^{-1}$ and
$c_i=\sqrt{1-p_i^2}$. The interior Jacobian $B=D_{\alpha_I}E$ is
tridiagonal, with
\begin{align}
 B_{ii}&=2+Rc_i(w_{i-1}+w_i),\\
 B_{i,i-1}&=Rw_{i-1}c_{i-1},\\
 B_{i,i+1}&=Rw_ic_{i+1},\\
 (D_{\alpha_0}E)_1&=Rw_0c_0,\\
 (D_{\alpha_m}E)_{m-1}&=Rw_{m-1}c_m.
 \label{eq:gr-jacobian}
\end{align}
Only entries with indices in $I$ belong to $B$; the last two entries
are boundary forcing. When $m=2$ both boundary terms are in the same
row. If $H=D^2S$ is the full stationary-action Hessian, then
\begin{equation}\label{eq:gr-hessian-factor}
 H_{II}=R\operatorname{diag}(c_i)_{i\in I}\,B.
\end{equation}
\end{lemma}
\begin{proof}
For an oriented physical flight $v_i=X_i/\tau_i$, the differential
of its direction is $D\vartheta_i=v_i^\perp\cdot DX_i/\tau_i$.
The departure and arrival normal components are respectively $c_i$
and $-c_{i+1}$. Since $n'(\alpha)=n(\alpha)^\perp$,
\[
 \partial_{\alpha_i}\vartheta_i=-Rc_iw_i,
 \qquad \partial_{\alpha_{i+1}}\vartheta_i=-Rc_{i+1}w_i.
\]
Applying these two identities to both flights in
\eqref{eq:gr-reflection} proves every entry, including both boundary
terms. Comparing with \eqref{eq:v109-second-variation} proves
\eqref{eq:gr-hessian-factor}. Thus a factor $c_i$ in the action Hessian
is a row factor in the gradient formulation; it is absent from the
reflection-angle Jacobian. This identity does not assert invertibility
of $H_{II}$ when a limiting $c_i$ is zero.
\end{proof}

\begin{theorem}[Column coercivity and a positive resolvent]
\label{thm:gr-resolvent}
Let $\mathcal S=\operatorname{diag}((-1)^i)_{i\in I}$ and
$A=\mathcal S B\mathcal S$. For every $m\ge2$, every original radius,
and every actual regular itinerary, $A^{-1}$ is entrywise nonnegative
and
\begin{equation}\label{eq:gr-l1}
 \|B^{-1}f\|_1\le\tfrac12\|f\|_1.
\end{equation}
More generally, for positive weights $v_i$ such that
$\lambda^{-1}\le v_{i+1}/v_i\le\lambda$, one has
\begin{equation}\label{eq:gr-weighted}
 \sum_{i\in I}v_i|(B^{-1}f)_i|
       \le\sum_{i\in I}v_i|f_i|.
\end{equation}
In particular,
\begin{equation}\label{eq:gr-green}
 \sum_{i\in I}\lambda^{|i-j|}(A^{-1})_{ij}\le1
 \quad(j\in I).
\end{equation}
The assertions extend algebraically to simultaneous values
$c_i=0$ by the same formulas. No lower incidence cutoff is required.
\end{theorem}
\begin{proof}
The off-diagonal entries of $A$ are nonpositive. Let
$D_{jj}=2+Rc_j(w_{j-1}+w_j)$ and write $A=(\mathrm{Id}-N)D$,
where $N$ has zero diagonal and nonnegative off-diagonal entries.
Its column sum is at most
\[
 \frac{Rc_j(w_{j-1}+w_j)}{2+Rc_j(w_{j-1}+w_j)}
 \le \frac{a_*}{1+a_*}=\frac{47}{53}<1.
\]
Consequently $A^{-1}=D^{-1}\sum_{k\ge0}N^k$ exists and is
nonnegative. This is an algebraic matrix series, not a probabilistic
replacement of the Lorentz gas. Missing interior neighbors at the two
ends only decrease the column sum.

Each column of $A$ has sum at least $2$:
$\boldsymbol1^{\mathsf T}A\ge2\boldsymbol1^{\mathsf T}$.
Multiplication on the right by the nonnegative inverse gives
$\boldsymbol1^{\mathsf T}A^{-1}\le\boldsymbol1^{\mathsf T}/2$.
Since $|B^{-1}f|\le A^{-1}|f|$, this proves \eqref{eq:gr-l1}.
The argument uses \emph{column}, not row, dominance. Row dominance
can fail when one incidence is much smaller than its neighbors.

For the weighted assertion, a column calculation gives
\[
 (v^{\mathsf T}A)_j
 \ge v_j\{2-(\lambda-1)Rc_j(w_{j-1}+w_j)\}
 \ge v_j\{2-2a_*(\lambda-1)\}=v_j.
\]
Here $2a_*(\lambda-1)=1$ exactly. Right multiplication by $A^{-1}$
proves \eqref{eq:gr-weighted}. Taking $v_i=\lambda^{|i-j|}$ proves
\eqref{eq:gr-green}. All denominators in this argument are at least
$2$. It therefore remains valid when any collection of incidences is
zero. No inverse $c_i$ has been introduced in taking that limit.
\end{proof}

\subsection{Complete higher endpoint derivatives}
For a vector $u=(u_0,\ldots,u_m)$ set
$V_i(u)=|u_i|+|u_{i+1}|$ and
$\|u\|_{1,v}=\sum_{i=0}^m v_i|u_i|$. We use
$v_i=\lambda^{d_i}$, $d_i=\min(i,m-i)$, including both endpoints.
For $r\ge1$ define the explicit constants
\begin{equation}\label{eq:gr-angle-constants}
 T_r=\sum_{k=1}^r {r\brace k}(k-1)!a_*^k,
 \qquad T_1=a_*,\quad T_2=a_*+a_*^2,
 \quad T_3=a_*+3a_*^2+2a_*^3,
\end{equation}
where ${r\brace k}$ counts partitions of $\{1,\ldots,r\}$ into $k$
nonempty blocks.

\begin{lemma}[Local angle tensors, with every chain-rule term]
\label{lem:gr-angle-tensors}
For each flight and each $r\ge1$,
\[
 |D^r\vartheta_i[u^1,\ldots,u^r]|
       \le T_r\prod_{h=1}^r V_i(u^h).
\]
For $r\ge2$, the full vector $E=(E_i)_{i\in I}$ consequently satisfies
\begin{align}
 \|D^rE[u^1,\ldots,u^r]\|_1
       &\le4T_r\prod_{h=1}^r\|u^h\|_1,\\
 \|D^rE[u^1,\ldots,u^r]\|_{1,v}
       &\le2(1+\lambda)T_r\prod_{h=1}^r\|u^h\|_{1,v}.
 \label{eq:gr-angle-norms}
\end{align}
\end{lemma}
\begin{proof}
Identify the plane with $\mathbb C$ and use a local logarithm of
$X_i\ne0$. Its $k$th derivative is
$(-1)^{k-1}(k-1)!X_i^{-k}$; taking the imaginary part gives the
direction derivative. Every derivative of
$X_i=O_{i+1}-O_i+R(e^{i\alpha_{i+1}}-e^{i\alpha_i})$
of order $b\ge1$ has absolute value at most
$R_+\prod_{h=1}^b V_i(u^h)$. The composition rule indexed by set
partitions gives exactly the sum in \eqref{eq:gr-angle-constants},
using $|X_i|\ge\tau_-$. This proves the first assertion without
omitting mixed derivatives.

Each flight appears in at most two interior reflection equations.
Also $V_i(u)\le\|u\|_1$ and $\sum_iV_i(u)\le2\|u\|_1$.
Summing one factor and bounding the others proves the first norm
inequality. For the second, the sum of the weights of the adjacent
reflection equations is at most $v_i+v_{i+1}$, and
\[
 (v_i+v_{i+1})V_i(u)
 \le(1+\lambda)(v_i|u_i|+v_{i+1}|u_{i+1}|).
\]
Summing over flights counts each vertex at most twice. The other
factors are bounded by their weighted norms, since $v_i\ge1$.
This proves the second inequality, including the boundary flights.
\end{proof}

Let $\mathfrak P_r$ denote the set of all set partitions of
$\{1,\ldots,r\}$. Define two finite positive recursions
\begin{align}
 C_1&=2+a_* =59/6,
 &C_r&=2\!\sum_{\substack{\pi\in\mathfrak P_r\\|\pi|\ge2}}
               T_{|\pi|}\prod_{Q\in\pi}C_{|Q|},\\
 M_1&=2+2\lambda a_* =56/3,
 &M_r&=2(1+\lambda)\!
       \sum_{\substack{\pi\in\mathfrak P_r\\|\pi|\ge2}}
               T_{|\pi|}\prod_{Q\in\pi}M_{|Q|},\\[-2pt]
 &&&\hspace{-18mm}r\ge2.\label{eq:gr-recursions}
\end{align}
The symbols $C_r$ and $M_r$ denote the jet constants in this section
and the next, not a redefinition of earlier flight-tensor constants.
Every block in these sums has size less than $r$, so the recursions
are not circular. In particular $C_2=2T_2C_1^2$ and
$C_3=2T_3C_1^3+6T_2C_1C_2$.

\begin{theorem}[All finite endpoint jets on actual physical branches]
\label{thm:gr-jets}
For every integer $r\ge1$, every $m\ge1$, and every actual regular
stationary lift $\alpha(q)$, $q=(\alpha_0,\alpha_m)$,
\begin{align}
 \|D_q^r\alpha[\eta^1,\ldots,\eta^r]\|_1
       &\le C_r\prod_{h=1}^r\|\eta^h\|_\infty,\\
 \sum_{i=0}^m\lambda^{d_i}
       |D_q^r\alpha_i[\eta^1,\ldots,\eta^r]|
       &\le M_r\prod_{h=1}^r\|\eta^h\|_\infty.
 \label{eq:gr-jets}
\end{align}
The constants depend on derivative order and the original geometric
interval only. They are independent of collision count, itinerary,
small positive incidence, return/displacement labels, or roof window.
\end{theorem}
\begin{proof}
For $r=1$, the interior equation is
$BJ_I=-D_{q}E\,\eta$. Its right side has norm at most
$a_*(|\eta_0|+|\eta_m|)$ and weighted norm at most
$\lambda a_*(|\eta_0|+|\eta_m|)$.
Theorem~\ref{thm:gr-resolvent}, followed by adding the two actual
endpoint derivatives, gives $C_1$ and $M_1$. This also covers $m=2$,
when the two boundary terms share a row. For $m=1$ only the two
endpoint derivatives occur.

For $r\ge2$ the endpoint derivatives of order $r$ vanish. Differentiate
$E(\alpha(q))=\mathrm{constant}$ in the $r$ prescribed directions.
Isolating the partition with one block gives
\[
 B(D_q^r\alpha)_I
 =-\!\sum_{\substack{\pi\in\mathfrak P_r\\|\pi|\ge2}}
 D^{|\pi|}E\bigl[D_q^{|Q|}\alpha[\eta^h:h\in Q]:Q\in\pi\bigr].
\]
This identity contains every mixed chain-rule term exactly once.
Induction, Lemma~\ref{lem:gr-angle-tensors}, and respectively the
$1/2$ and $1$ inverse bounds yield precisely
\eqref{eq:gr-recursions}. Local smoothness follows already from the
ordinary implicit function theorem and invertibility of $B$; the
induction establishes uniform bounds for its derivatives. No assertion
of a common physical continuation ball is used.
\end{proof}

\begin{corollary}[Fixed-itinerary grazing limits, without a fictitious event]
\label{cor:gr-grazing-limit}
Fix $m$ and a sequence of actual regular trajectories with the same
centers and convergent collision angles and radius. Suppose the limit
is in the original radius interval. Any number of interior incidences
may tend to zero. The local reflection equations at the limiting
configuration have a unique smooth formal solution in the endpoint
angles; every finite endpoint jet of the original physical branch has
that formal limit. The bounds in \eqref{eq:gr-jets} persist.
This does not make the formal continuation an admissible billiard
trajectory beyond its actual first boundary.
\end{corollary}
\begin{proof}
Distinct successive disk centers and $R\le R_+$ keep every limiting
flight length at least $\tau_-$. Local direction logarithms are smooth
there. The limiting Jacobian is invertible by
Theorem~\ref{thm:gr-resolvent}, including simultaneous zero incidences.
The implicit function theorem at that configuration supplies the
formal reflection solution, with the radius also treated as a smooth
parameter. Local uniqueness identifies every
sufficiently close member of the given physical sequence with it.
Its finite jets converge, and their uniform bounds pass to the limit.
An extra disk can obstruct the formal solution, and the prescribed
section or return event can end at the boundary. Neither obstruction
is discarded; no collision-map derivative at tangency is asserted.
\end{proof}
